Backgammon combines visible board positions with stochastic dice outcomes. Each player moves 15 checkers around 24 points toward a home board and then bears them off. A point occupied by two or more opposing checkers is blocked. Landing on a lone opposing checker, or blot, sends it to the bar. Under the implemented standard movement rules, bar checkers must re-enter before other checkers move. Legal complete turns must use the maximum possible number of dice, with the higher die required when only one of two different dice can be played. Doubles allow four uses of the rolled value.

A decision occurs after the dice are known. Flat Monte Carlo search evaluates each legal candidate by sampling future play and averaging its outcomes. Unlike Monte Carlo Tree Search, it does not maintain and progressively expand a tree or use a tree-selection rule. The survey by Browne et al. provides broader context for simulation-based tree-search methods \cite{browne2012}; the present implementation uses the simpler flat approach.

Full-game simulations were implemented initially, but their cost motivated a configurable depth limit. A depth of four means four individual player turns after the candidate move: opponent, original player, opponent, original player. A forced pass consumes one turn. If a simulation has not ended at the limit, a position heuristic supplies its score. This is a utility estimate, not a calibrated winning probability.
